\subsection{Cross-Domain Patterns: Five Principles for Behavioral Intervention}
\label{subsec:cross-domain-patterns}

%%%%%%%%%%%%%%%%%%%%%%%%%%%%%%%%%%%%%%%%%%%%%%%%%%%%%%%%%%%%%%%%%%%%%%%%%%%%%%%
% 10.10: Systematic analysis of patterns emerging across all nine domains
%%%%%%%%%%%%%%%%%%%%%%%%%%%%%%%%%%%%%%%%%%%%%%%%%%%%%%%%%%%%%%%%%%%%%%%%%%%%%%%

The nine application domains (Section~\ref{subsec:applications}) revealed 
recurring patterns. This section systematically develops these patterns,
connecting them to the theoretical framework (Sections~\ref{subsec:three-categories}--\ref{subsec:aggregate-welfare})
and deriving practical principles for behavioral intervention.

\subsubsection{Overview: Five Cross-Domain Patterns}

\begin{table}[htbp]
\centering
\caption{Five Cross-Domain Patterns}
\label{tab:five-patterns}
\renewcommand{\arraystretch}{1.4}
\begin{tabular}{@{}clp{7cm}@{}}
\toprule
\textbf{\#} & \textbf{Pattern} & \textbf{Core Statement} \\
\midrule
10.10.1 & The INU-Only Trap & Interventions targeting only individual utility systematically fail \\
10.10.2 & The Underestimated IDN & Identity utility is the most neglected yet often most powerful driver \\
10.10.3 & The $\gamma$-Sustainability Principle & Change sustainability depends on complementarity structure \\
10.10.4 & The $\Psi$-First Sequence & Context must be addressed before outcomes \\
10.10.5 & The Discounting Rationality & ``Irrational'' behavior is often rational given time structure \\
\bottomrule
\end{tabular}
\end{table}

%%%%%%%%%%%%%%%%%%%%%%%%%%%%%%%%%%%%%%%%%%%%%%%%%%%%%%%%%%%%%%%%%%%%%%%%%%%%%%%
\subsubsection{Pattern 10.10.1: The INU-Only Trap}
\label{subsubsec:inu-only-trap}
%%%%%%%%%%%%%%%%%%%%%%%%%%%%%%%%%%%%%%%%%%%%%%%%%%%%%%%%%%%%%%%%%%%%%%%%%%%%%%%

\paragraph{The Pattern.}
Interventions that target only Individual Utility ($U^{INU}$) systematically 
underperform or fail, even when the INU case is strong.

\paragraph{Theoretical Foundation.}
From the aggregate welfare function (Section~\ref{subsec:aggregate-welfare}):

\begin{equation}
Q = \omega_{INU} \cdot U^{INU} + \omega_{KNU} \cdot U^{KNU} + \omega_{IDN} \cdot U^{IDN} + \sum_{i<j} \gamma_{ij} \cdot U^i \cdot U^j
\end{equation}

An INU-only intervention changes only $\Delta U^{INU}$:

\begin{equation}
\Delta Q_{INU-only} = \omega_{INU} \cdot \Delta U^{INU} + \gamma_{INU,KNU} \cdot U^{KNU} \cdot \Delta U^{INU} + \gamma_{INU,IDN} \cdot U^{IDN} \cdot \Delta U^{INU}
\end{equation}

If $\gamma_{INU,KNU} < 0$ or $\gamma_{INU,IDN} < 0$, the complementarity terms 
\emph{offset} the direct INU gain. With typical weights ($\omega_{INU} \approx 0.4$),
even a large $\Delta U^{INU}$ may not dominate negative complementarity effects.

\paragraph{Evidence Across Domains.}

\begin{table}[htbp]
\centering
\caption{The INU-Only Trap Across Domains}
\label{tab:inu-trap-evidence}
\renewcommand{\arraystretch}{1.3}
\begin{tabular}{@{}llll@{}}
\toprule
\textbf{Domain} & \textbf{INU-Only Intervention} & \textbf{Why It Fails} & \textbf{Ref.} \\
\midrule
Labor & Retraining subsidies & Ignores IDN (identity loss) & 10.9.1 \\
Health & Health information & Ignores IDN (``smoker'' identity) & 10.9.2 \\
Finance & Financial incentives & Ignores E, S, IDN functions of spending & 10.9.3 \\
Sustainability & Carbon tax alone & Ignores KNU (collective action) & 10.9.4 \\
Digital & Efficiency arguments & Ignores IDN (``I'm obsolete'') & 10.9.5 \\
Education & Tuition subsidies & Ignores IDN (adult learner stigma) & 10.9.7 \\
\bottomrule
\end{tabular}
\end{table}

\paragraph{Why Do Policymakers Fall Into This Trap?}

\begin{enumerate}
\item \textbf{Economic training:} Standard economics focuses on INU (homo economicus)
\item \textbf{Measurability:} $U^{INU}_F$ is easy to measure; $U^{IDN}$ is not
\item \textbf{Policy tools:} Taxes and subsidies directly affect INU
\item \textbf{Projection:} Policymakers may weight INU more than target populations
\end{enumerate}

\paragraph{The Mathematical Condition for INU-Only Success.}
INU-only interventions succeed only when:

\begin{equation}
\omega_{INU} \cdot \Delta U^{INU} > |\gamma_{INU,KNU} \cdot U^{KNU} \cdot \Delta U^{INU}| + |\gamma_{INU,IDN} \cdot U^{IDN} \cdot \Delta U^{INU}|
\end{equation}

This requires either:
\begin{itemize}
\item Very high $\omega_{INU}$ (rare outside pure financial decisions)
\item Non-negative $\gamma_{INU,KNU}$ and $\gamma_{INU,IDN}$ (aligned categories)
\item Very large $\Delta U^{INU}$ (expensive)
\end{itemize}

\paragraph{Practical Principle.}

\begin{tcolorbox}[colback=blue!5!white, colframe=blue!75!black, title=Principle 1: Beyond INU]
\textbf{Before designing any intervention, ask:}
\begin{enumerate}
\item What is the KNU impact? (collective welfare, fairness)
\item What is the IDN impact? (identity confirmation or threat)
\item Are the complementarities positive or negative?
\end{enumerate}
\textbf{Rule of thumb:} If the intervention only has an INU story, it will 
likely underperform. Add KNU and IDN components.
\end{tcolorbox}

%%%%%%%%%%%%%%%%%%%%%%%%%%%%%%%%%%%%%%%%%%%%%%%%%%%%%%%%%%%%%%%%%%%%%%%%%%%%%%%
\subsubsection{Pattern 10.10.2: The Underestimated IDN}
\label{subsubsec:underestimated-idn}
%%%%%%%%%%%%%%%%%%%%%%%%%%%%%%%%%%%%%%%%%%%%%%%%%%%%%%%%%%%%%%%%%%%%%%%%%%%%%%%

\paragraph{The Pattern.}
Identity Utility ($U^{IDN}$) is systematically underestimated in intervention 
design, yet is often the decisive factor in behavioral outcomes.

\paragraph{Theoretical Foundation.}
From Section~\ref{subsec:three-categories}, IDN answers: ``Who am I?''

\begin{equation}
U^{IDN}_d = f(\text{match between action and self-concept})
\end{equation}

Three properties make IDN uniquely powerful:

\begin{enumerate}
\item \textbf{Immediacy:} IDN effects are felt instantly ($t=0$), not discounted
\item \textbf{Stickiness:} IDN reference points change slowly (Section~\ref{subsec:reference-points})
\item \textbf{Non-fungibility:} IDN cannot be compensated by INU gains
\end{enumerate}

\paragraph{The Immediacy Advantage.}
While INU health benefits are discounted ($\delta^3 \approx 0.6$), 
IDN effects are immediate:

\begin{align}
U^{INU}_{P,3}(\text{discounted}) &= \delta^3 \cdot U^{INU}_{P,3} \approx 0.6 \cdot U^{INU}_{P,3} \\
U^{IDN}_{E,0}(\text{not discounted}) &= 1.0 \cdot U^{IDN}_{E,0}
\end{align}

An identity confirmation \emph{now} has higher present value than a health 
benefit \emph{later}.

\paragraph{The Stickiness Problem.}
IDN reference points have the lowest adaptation rate (Section~\ref{subsec:reference-points}):

\begin{equation}
\eta^{IDN} \approx 0.1-0.3 \quad \text{vs.} \quad \eta^{INU} \approx 0.3-0.5
\end{equation}

This means:
\begin{itemize}
\item A ``coal miner'' identity persists long after the job is gone
\item A ``smoker'' identity persists long after quitting attempts
\item A ``non-athletic'' identity blocks exercise adoption for years
\end{itemize}

\paragraph{The Non-Fungibility Property.}
From Section~\ref{subsec:inter-category-complementarities}, when $\gamma_{INU,IDN} < 0$:

\begin{equation}
\text{More } U^{INU} \text{ with violated IDN} \Rightarrow \text{Lower total } Q
\end{equation}

You cannot simply ``pay people'' to violate their identity---the money 
makes it \emph{worse} (``selling out'').

\paragraph{Evidence Across Domains.}

\begin{table}[htbp]
\centering
\caption{IDN as Decisive Factor Across Domains}
\label{tab:idn-evidence}
\renewcommand{\arraystretch}{1.3}
\begin{tabular}{@{}lll@{}}
\toprule
\textbf{Domain} & \textbf{IDN Factor} & \textbf{Effect} \\
\midrule
Labor & ``I'm a miner'' & Blocks retraining despite +\euro{20k} \\
Health & ``I'm a smoker'' & Blocks cessation despite health info \\
Sustainability & ``I'm an environmentalist'' & Enables sacrifice despite INU cost \\
Digital & ``I'm an expert'' & Creates resistance to new systems \\
Education & ``I'm not a student'' & Blocks adult learning \\
Marketing & Brand identity & Creates loyalty beyond product quality \\
\bottomrule
\end{tabular}
\end{table}

\paragraph{Why Is IDN Underestimated?}

\begin{enumerate}
\item \textbf{Invisibility:} Identity is internal, not directly observable
\item \textbf{Measurement difficulty:} No standard IDN metrics exist
\item \textbf{Rationalist bias:} Assumes people respond to ``objective'' incentives
\item \textbf{Projection error:} Experts may have weaker occupational identity
\item \textbf{Taboo:} Discussing identity feels ``soft'' or ``unscientific''
\end{enumerate}

\paragraph{When IDN Dominates.}
IDN becomes the decisive factor when:

\begin{itemize}
\item The behavior is publicly visible (identity signaling)
\item The behavior has symbolic meaning (``what does this say about me?'')
\item Core roles are threatened (career, family, community position)
\item The person has strong, crystallized identity
\item The action would violate deeply held values
\end{itemize}

\paragraph{Practical Principle.}

\begin{tcolorbox}[colback=blue!5!white, colframe=blue!75!black, title=Principle 2: Identity First]
\textbf{For any behavior change intervention:}
\begin{enumerate}
\item \textbf{Diagnose:} What identity is currently attached to the behavior?
\item \textbf{Assess:} Would the new behavior violate or confirm this identity?
\item \textbf{Design:} Create an ``identity bridge'' if violation is likely
\item \textbf{Frame:} ``Become a new version of yourself'' rather than ``give up who you are''
\end{enumerate}
\textbf{The identity bridge:} Find elements of the old identity that connect 
to the new behavior. Coal miner → ``Builder of the energy future.''
\end{tcolorbox}

%%%%%%%%%%%%%%%%%%%%%%%%%%%%%%%%%%%%%%%%%%%%%%%%%%%%%%%%%%%%%%%%%%%%%%%%%%%%%%%
\subsubsection{Pattern 10.10.3: The $\gamma$-Sustainability Principle}
\label{subsubsec:gamma-sustainability}
%%%%%%%%%%%%%%%%%%%%%%%%%%%%%%%%%%%%%%%%%%%%%%%%%%%%%%%%%%%%%%%%%%%%%%%%%%%%%%%

\paragraph{The Pattern.}
The sustainability of behavioral change depends critically on the 
complementarity structure ($\gamma_{ij}$). Changes that disrupt positive 
complementarities tend to reverse; changes that create positive 
complementarities tend to persist.

\paragraph{Theoretical Foundation.}
From Section~\ref{subsec:inter-category-complementarities}, the complementarities 
determine whether utility categories reinforce or undermine each other:

\begin{table}[htbp]
\centering
\caption{Complementarity and Change Sustainability}
\label{tab:gamma-sustainability}
\renewcommand{\arraystretch}{1.3}
\begin{tabular}{@{}lll@{}}
\toprule
\textbf{$\gamma$-Structure} & \textbf{Change Dynamic} & \textbf{Outcome} \\
\midrule
All $\gamma > 0$ (coherent) & Self-reinforcing & Change persists \\
Mixed $\gamma$ & Partial tension & Change vulnerable \\
All $\gamma < 0$ (conflicted) & Self-undermining & Change reverses \\
\bottomrule
\end{tabular}
\end{table}

\paragraph{The Mechanism: Reinforcement vs. Erosion.}
When complementarities are positive:
\begin{equation}
\Delta U^{INU} > 0 \Rightarrow \frac{\partial U^{KNU}}{\partial U^{INU}} > 0 \Rightarrow \Delta U^{KNU} > 0 \Rightarrow \ldots
\end{equation}

Gains in one category amplify gains in others---a virtuous cycle.

When complementarities are negative:
\begin{equation}
\Delta U^{INU} > 0 \Rightarrow \frac{\partial U^{KNU}}{\partial U^{INU}} < 0 \Rightarrow \Delta U^{KNU} < 0 \Rightarrow \text{pressure to reverse}
\end{equation}

The KNU loss creates pressure to undo the INU gain---relapse.

\paragraph{The Relapse Prediction.}
From the Journey model (Chapter~\ref{sec:journey}, preview):

\begin{equation}
P(\text{Relapse}) = f(\gamma_{structure})
\end{equation}

\begin{itemize}
\item Positive $\gamma$: New behavior satisfies multiple categories → low relapse
\item Negative $\gamma$: New behavior creates tension → high relapse
\end{itemize}

\paragraph{Evidence Across Domains.}

\begin{table}[htbp]
\centering
\caption{$\gamma$-Structure and Sustainability}
\label{tab:gamma-evidence}
\renewcommand{\arraystretch}{1.3}
\begin{tabular}{@{}llll@{}}
\toprule
\textbf{Domain} & \textbf{Behavior Change} & \textbf{$\gamma$-Structure} & \textbf{Sustainability} \\
\midrule
Health & Solo quitting & $\gamma_{INU,KNU} < 0$ (leaves smoking group) & Low (high relapse) \\
Health & Group quitting & $\gamma_{INU,KNU} > 0$ (new support group) & High \\
Labor & Individual retraining & All $\gamma < 0$ (leaves community, identity) & Very low \\
Labor & Community transition & $\gamma_{KNU,IDN} > 0$ (``we transform together'') & Higher \\
Culture & Top-down mandate & $\gamma_{INU,KNU} < 0$ (imposed, not owned) & Low \\
Culture & Participatory change & All $\gamma > 0$ (aligned interests) & High \\
\bottomrule
\end{tabular}
\end{table}

\paragraph{The Coal Miner Test.}
The coal miner transition (Section~\ref{subsubsec:labor-transitions}) 
perfectly illustrates this pattern:

\begin{itemize}
\item \textbf{Status quo:} All $\gamma > 0$ (coherent miner identity)
\item \textbf{After individual transition:} All $\gamma < 0$ (conflicted coder)
\item \textbf{Prediction:} Even if forced to transition, high ``relapse'' risk 
(leaves coding, returns to mining community, or experiences severe distress)
\end{itemize}

\paragraph{Designing for Positive $\gamma$.}
Sustainable interventions must either:

\begin{enumerate}
\item \textbf{Preserve existing positive $\gamma$:} Don't disrupt what works
\item \textbf{Transform $\gamma$ structure:} Build new complementarities
\item \textbf{Move the whole system:} Collective transitions preserve KNU
\end{enumerate}

\paragraph{Practical Principle.}

\begin{tcolorbox}[colback=blue!5!white, colframe=blue!75!black, title=Principle 3: Complementarity Design]
\textbf{For sustainable behavioral change:}
\begin{enumerate}
\item \textbf{Map current $\gamma$:} What complementarities exist now?
\item \textbf{Project future $\gamma$:} What happens to complementarities after change?
\item \textbf{Design for coherence:} Create conditions where new behavior 
aligns INU, KNU, and IDN
\end{enumerate}
\textbf{Warning signs:} If the intervention requires sustained willpower 
against multiple category pressures, it will likely fail. Sustainable 
change feels \emph{natural} because all categories align.
\end{tcolorbox}

%%%%%%%%%%%%%%%%%%%%%%%%%%%%%%%%%%%%%%%%%%%%%%%%%%%%%%%%%%%%%%%%%%%%%%%%%%%%%%%
\subsubsection{Pattern 10.10.4: The $\Psi$-First Sequence}
\label{subsubsec:psi-first}
%%%%%%%%%%%%%%%%%%%%%%%%%%%%%%%%%%%%%%%%%%%%%%%%%%%%%%%%%%%%%%%%%%%%%%%%%%%%%%%

\paragraph{The Pattern.}
Interventions that improve outcomes ($x$) without first addressing context 
($\Psi$) yield suboptimal welfare gains. The sequence matters: $\Psi$ first, 
then $\gamma$, then $x$.

\paragraph{Theoretical Foundation.}
From Section~\ref{subsec:context-modulation}, all welfare parameters depend on context:

\begin{align}
\delta(\Psi_T) &: \text{Discounting depends on stability} \\
\lambda(\Psi_I) &: \text{Loss aversion depends on institutions} \\
\gamma_{ij}(\Psi_S) &: \text{Complementarities depend on trust} \\
r^i_d(\Psi) &: \text{Reference points depend on context}
\end{align}

The same outcome $x$ produces different welfare in different contexts:

\begin{equation}
Q(x, \Psi_{good}) >> Q(x, \Psi_{bad})
\end{equation}

\paragraph{The K-Q-$\Psi$ Chain.}
From Section~\ref{subsec:aggregate-welfare}:

\begin{equation}
\Psi \longrightarrow K \longrightarrow Q
\end{equation}

Context enables coherence; coherence enables welfare realization.
Without adequate $\Psi$, even good outcomes fail to translate into welfare.

\paragraph{Evidence Across Domains.}

\begin{table}[htbp]
\centering
\caption{$\Psi$-First Sequence Evidence}
\label{tab:psi-first-evidence}
\renewcommand{\arraystretch}{1.3}
\begin{tabular}{@{}lp{5cm}p{5cm}@{}}
\toprule
\textbf{Domain} & \textbf{Outcome-First Failure} & \textbf{$\Psi$-First Success} \\
\midrule
Labor & Retraining without community support & Community transition with infrastructure \\
Finance & Savings accounts without financial literacy & Education first, then products \\
Sustainability & Carbon tax without trust & Build trust, then introduce policy \\
Culture & New values statement without leadership change & Model behavior, then formalize \\
Policy & Regulation without institutional capacity & Build institutions, then regulate \\
\bottomrule
\end{tabular}
\end{table}

\paragraph{Why Does Outcome-First Fail?}
Several mechanisms:

\begin{enumerate}
\item \textbf{High discounting:} In unstable contexts ($\Psi_T$ low), future 
benefits are heavily discounted → present costs dominate
\item \textbf{Negative complementarities:} In low-trust contexts ($\Psi_S$ low), 
individual gains create suspicion → $\gamma_{INU,KNU} < 0$
\item \textbf{Weak institutions:} Without $\Psi_I$, gains cannot be protected → 
people don't invest in change
\item \textbf{Reference point instability:} Without $\Psi_T$, reference points 
fluctuate → no stable baseline for Gain/Pain calculation
\end{enumerate}

\paragraph{The Sequencing Principle.}
Optimal intervention sequence:

\begin{enumerate}
\item \textbf{First: $\Psi$} (Context)
\begin{itemize}
\item Build institutions ($\Psi_I$)
\item Establish trust ($\Psi_S$)
\item Create stability ($\Psi_T$)
\end{itemize}

\item \textbf{Second: $\gamma$} (Complementarities)
\begin{itemize}
\item Align incentives so $\gamma_{INU,KNU} > 0$
\item Build identity bridges so $\gamma_{INU,IDN} > 0$
\item Strengthen group identity so $\gamma_{KNU,IDN} > 0$
\end{itemize}

\item \textbf{Third: $x$} (Outcomes)
\begin{itemize}
\item Now improve outcomes
\item They will translate efficiently to welfare
\item Change will be sustainable
\end{itemize}
\end{enumerate}

\paragraph{The Impatience Problem.}
Why do interventions skip the sequence?

\begin{itemize}
\item \textbf{Political pressure:} ``We need results now''
\item \textbf{Measurability:} Outcomes are visible; $\Psi$ is not
\item \textbf{Time horizons:} $\Psi$-building takes years
\item \textbf{Attribution:} Hard to credit $\Psi$-work for later success
\end{itemize}

\paragraph{Practical Principle.}

\begin{tcolorbox}[colback=blue!5!white, colframe=blue!75!black, title=Principle 4: Sequence Matters]
\textbf{The optimal intervention sequence:}

\begin{center}
\textbf{$\Psi$ (Context)} $\longrightarrow$ \textbf{$\gamma$ (Alignment)} $\longrightarrow$ \textbf{$x$ (Outcomes)}
\end{center}

\textbf{Diagnostic questions:}
\begin{enumerate}
\item Is the context adequate? ($\Psi_I$, $\Psi_S$, $\Psi_T$ sufficient?)
\item Are categories aligned? (All $\gamma > 0$?)
\item Only then: Will improving $x$ increase $Q$?
\end{enumerate}

\textbf{Rule:} If context is poor, invest in $\Psi$ first. The same resources 
spent on outcome improvement in good context yield 2-3x the welfare gain.
\end{tcolorbox}

%%%%%%%%%%%%%%%%%%%%%%%%%%%%%%%%%%%%%%%%%%%%%%%%%%%%%%%%%%%%%%%%%%%%%%%%%%%%%%%
\subsubsection{Pattern 10.10.5: The Discounting Rationality}
\label{subsubsec:discounting-rationality}
%%%%%%%%%%%%%%%%%%%%%%%%%%%%%%%%%%%%%%%%%%%%%%%%%%%%%%%%%%%%%%%%%%%%%%%%%%%%%%%

\paragraph{The Pattern.}
Behaviors labeled ``irrational''---undersaving, unhealthy choices, 
short-term thinking---are often \emph{rational} responses to time structure
and context. Heavy discounting is an adaptive response to uncertainty.

\paragraph{Theoretical Foundation.}
From Section~\ref{subsec:calculation-logic}, time discounting enters the 
utility calculation:

\begin{equation}
U^i_{d,t} = (\delta^i_d)^t \cdot \left[ G^i_{d,t} - \lambda^i_d \cdot P^i_{d,t} \right]
\end{equation}

With $\delta = 0.85$:
\begin{align}
t=0: \quad & \delta^0 = 1.00 \quad \text{(full weight)} \\
t=1: \quad & \delta^1 = 0.85 \quad \text{(85\% weight)} \\
t=2: \quad & \delta^2 = 0.72 \quad \text{(72\% weight)} \\
t=3: \quad & \delta^3 = 0.61 \quad \text{(61\% weight)}
\end{align}

A long-term benefit must be ~1.6x larger than a short-term cost to be preferred.

\paragraph{Context-Dependent Discounting.}
From Section~\ref{subsec:context-modulation}:

\begin{equation}
\delta(\Psi) = \bar{\delta} + \gamma_{\delta,T} \cdot \Psi_T
\end{equation}

In unstable contexts ($\Psi_T$ low):
\begin{itemize}
\item Future is genuinely uncertain
\item Rational to discount heavily
\item ``Short-term thinking'' is adaptive
\end{itemize}

\paragraph{The Rationality Reframe.}
What appears ``irrational'':

\begin{table}[htbp]
\centering
\caption{Reframing ``Irrational'' Behaviors}
\label{tab:rationality-reframe}
\renewcommand{\arraystretch}{1.3}
\begin{tabular}{@{}lp{5cm}p{5cm}@{}}
\toprule
\textbf{``Irrational'' Behavior} & \textbf{Standard View} & \textbf{EBF View} \\
\midrule
Undersaving & Failure of self-control & Rational given uncertainty ($\Psi_T$ low) \\
Unhealthy choices & Ignorance or weakness & Heavy discounting of $U^P_{t=3}$ \\
Present consumption & Myopia & Immediate $U^E$, $U^S$, $U^{IDN}$ gains \\
Rejecting retraining & Irrationality & Protecting $U^{IDN}$ and $U^{KNU}$ \\
\bottomrule
\end{tabular}
\end{table}

\paragraph{Evidence Across Domains.}

\begin{table}[htbp]
\centering
\caption{Discounting Rationality Across Domains}
\label{tab:discounting-evidence}
\renewcommand{\arraystretch}{1.3}
\begin{tabular}{@{}llp{5cm}@{}}
\toprule
\textbf{Domain} & \textbf{Behavior} & \textbf{Discount-Rational Explanation} \\
\midrule
Health & Smoking despite knowledge & $U^E_{t=0}$ (relief) >> $\delta^3 \cdot P^P_{t=3}$ \\
Finance & Undersaving & $\Psi_T$ low → future uncertain → discount heavily \\
Education & Not investing in learning & Uncertain payback, certain immediate cost \\
Sustainability & Ignoring climate & $\delta^{30} \approx 0$ for 30-year horizon \\
\bottomrule
\end{tabular}
\end{table}

\paragraph{The Poverty Trap Mechanism.}
Low $\Psi$ creates a rational discounting trap:

\begin{equation}
\Psi_T \downarrow \quad \Rightarrow \quad \delta \downarrow \quad \Rightarrow \quad \text{Don't invest in future} \quad \Rightarrow \quad \Psi_T \downarrow \quad \Rightarrow \quad \ldots
\end{equation}

Poverty and instability create \emph{rational} short-term orientation,
which perpetuates poverty and instability.

\paragraph{Breaking the Cycle.}
Two approaches:

\begin{enumerate}
\item \textbf{Increase $\Psi_T$:} Make the future more certain
\begin{itemize}
\item Social insurance reduces downside risk
\item Stable institutions make investment safer
\item Long-term contracts extend time horizons
\end{itemize}

\item \textbf{Front-load benefits:} Restructure the time pattern
\begin{itemize}
\item Immediate rewards for long-term behaviors
\item Quick wins before delayed payoffs
\item Reduce upfront costs (subsidies, defaults)
\end{itemize}
\end{enumerate}

\paragraph{Practical Principle.}

\begin{tcolorbox}[colback=blue!5!white, colframe=blue!75!black, title=Principle 5: Respect the Discount Rate]
\textbf{``Irrational'' behavior is often rational given time structure.}

\textbf{Before labeling behavior irrational, ask:}
\begin{enumerate}
\item What is the person's effective $\delta$? (Context-dependent!)
\item What is the time pattern of costs and benefits?
\item Is ``short-term thinking'' actually adaptive in their context?
\end{enumerate}

\textbf{Intervention approaches:}
\begin{itemize}
\item \textbf{Increase $\delta$:} Stabilize context, reduce uncertainty
\item \textbf{Restructure timing:} Front-load benefits, delay costs
\item \textbf{Create commitment:} Lock in future behavior when present self is motivated
\end{itemize}

\textbf{Warning:} Calling behavior ``irrational'' and demanding change 
without addressing the discount-rate drivers will fail.
\end{tcolorbox}

%%%%%%%%%%%%%%%%%%%%%%%%%%%%%%%%%%%%%%%%%%%%%%%%%%%%%%%%%%%%%%%%%%%%%%%%%%%%%%%
\subsubsection{Integration: The Five Principles in Practice}
\label{subsubsec:five-principles-integration}
%%%%%%%%%%%%%%%%%%%%%%%%%%%%%%%%%%%%%%%%%%%%%%%%%%%%%%%%%%%%%%%%%%%%%%%%%%%%%%%

\paragraph{The Principles Work Together.}
The five patterns are not independent---they form an integrated framework 
for behavioral intervention:

\begin{table}[htbp]
\centering
\caption{The Five Principles: Integration}
\label{tab:principles-integration}
\renewcommand{\arraystretch}{1.4}
\begin{tabular}{@{}cp{10cm}@{}}
\toprule
\textbf{\#} & \textbf{Principle and Integration} \\
\midrule
1 & \textbf{Beyond INU:} Don't target only individual utility \\
& $\rightarrow$ leads to Principle 2 (check IDN) and Principle 3 (check $\gamma$) \\
\midrule
2 & \textbf{Identity First:} Diagnose and design for identity \\
& $\rightarrow$ requires Principle 3 (build positive $\gamma_{INU,IDN}$) \\
\midrule
3 & \textbf{Complementarity Design:} Create self-reinforcing structures \\
& $\rightarrow$ requires Principle 4 ($\Psi$ enables positive $\gamma$) \\
\midrule
4 & \textbf{$\Psi$-First Sequence:} Build context before outcomes \\
& $\rightarrow$ enables Principle 5 (stable $\Psi_T$ increases $\delta$) \\
\midrule
5 & \textbf{Respect Discounting:} Address time structure \\
& $\rightarrow$ circles back to Principle 1 (immediate IDN/KNU vs. delayed INU) \\
\bottomrule
\end{tabular}
\end{table}

\paragraph{The Complete Diagnostic.}
For any behavioral intervention:

\begin{enumerate}
\item \textbf{Context check:} Is $\Psi$ adequate? (Principle 4)
\begin{itemize}
\item If no: Invest in $\Psi$ first
\end{itemize}

\item \textbf{Discounting check:} What is effective $\delta$? (Principle 5)
\begin{itemize}
\item If low: Address uncertainty or restructure timing
\end{itemize}

\item \textbf{Category check:} Which categories are affected? (Principle 1)
\begin{itemize}
\item If INU-only: Add KNU and IDN components
\end{itemize}

\item \textbf{Identity check:} What identity is at stake? (Principle 2)
\begin{itemize}
\item If threatened: Build identity bridge
\end{itemize}

\item \textbf{Complementarity check:} What is $\gamma$-structure? (Principle 3)
\begin{itemize}
\item If negative: Redesign for coherence
\end{itemize}
\end{enumerate}

\paragraph{Application: The Complete Coal Miner Intervention.}
Applying all five principles to labor transition:

\begin{table}[htbp]
\centering
\caption{Coal Miner Transition: Five-Principle Design}
\label{tab:miner-five-principles}
\renewcommand{\arraystretch}{1.3}
\begin{tabular}{@{}lp{8cm}@{}}
\toprule
\textbf{Principle} & \textbf{Application} \\
\midrule
1. Beyond INU & Not just retraining subsidy; include community support (KNU) and identity work (IDN) \\
2. Identity First & ``Energy transition leader'' identity bridge \\
3. $\gamma$-Design & Community transitions together (preserves $\gamma_{KNU,IDN}$) \\
4. $\Psi$-First & Build infrastructure and safety nets before expecting transition \\
5. Discounting & Front-load benefits; guaranteed income during transition \\
\bottomrule
\end{tabular}
\end{table}

%%%%%%%%%%%%%%%%%%%%%%%%%%%%%%%%%%%%%%%%%%%%%%%%%%%%%%%%%%%%%%%%%%%%%%%%%%%%%%%
\subsubsection{Summary: Five Cross-Domain Patterns}
%%%%%%%%%%%%%%%%%%%%%%%%%%%%%%%%%%%%%%%%%%%%%%%%%%%%%%%%%%%%%%%%%%%%%%%%%%%%%%%

\begin{tcolorbox}[colback=gray!5!white, colframe=gray!75!black, title=Section 10.10 Summary: Five Principles for Behavioral Intervention]

\textbf{Pattern 1: The INU-Only Trap}
\begin{itemize}
\item Interventions targeting only $U^{INU}$ systematically fail
\item Must address KNU (collective) and IDN (identity) components
\end{itemize}

\textbf{Pattern 2: The Underestimated IDN}
\begin{itemize}
\item Identity utility is immediate, sticky, and non-fungible
\item Often the decisive factor despite being invisible
\item Design identity bridges for behavior change
\end{itemize}

\textbf{Pattern 3: The $\gamma$-Sustainability Principle}
\begin{itemize}
\item Change sustainability depends on complementarity structure
\item Positive $\gamma$: self-reinforcing, persistent
\item Negative $\gamma$: self-undermining, relapse
\end{itemize}

\textbf{Pattern 4: The $\Psi$-First Sequence}
\begin{itemize}
\item Context before outcomes: $\Psi \rightarrow \gamma \rightarrow x$
\item Good outcomes in bad context yield poor welfare
\item Invest in institutions, trust, stability first
\end{itemize}

\textbf{Pattern 5: The Discounting Rationality}
\begin{itemize}
\item ``Irrational'' behavior often rational given time structure
\item Heavy discounting is adaptive in unstable contexts
\item Address $\Psi_T$ or restructure timing
\end{itemize}

\textbf{Integration:}
The five principles form a complete diagnostic and design framework.
Apply systematically to any behavioral intervention challenge.

\end{tcolorbox}

%%%%%%%%%%%%%%%%%%%%%%%%%%%%%%%%%%%%%%%%%%%%%%%%%%%%%%%%%%%%%%%%%%%%%%%%%%%%%%%
% END OF SECTION 10.10 AND CHAPTER 10
%%%%%%%%%%%%%%%%%%%%%%%%%%%%%%%%%%%%%%%%%%%%%%%%%%%%%%%%%%%%%%%%%%%%%%%%%%%%%%%
